\chapter{Barrick Mining Corporation and Valuation Perimeter}
\chapterstatus{Legal identity and listing continuity are verified from Barrick's
official releases. The gold operating perimeter, variable definitions, units,
data window and forecast frequency follow the Team~4 source study
\cite{team4}; ownership shares, accounting treatment and 2025 structural events
are documented in Barrick's 2025 Annual Report \cite{barrick-ar-2025}.
A separate copper extension is identified explicitly. The financial perimeter
and equity bridge remain controlled inputs requiring reconciliation.}

\section{Corporate identity and the choice of case study}
\label{sec:company-identity}

The project retained its historical working name, ``Barrick Gold'', but the
legal issuer is now \emph{Barrick Mining Corporation}. Shareholders approved
the change from Barrick Gold Corporation on 6 May 2025. Trading under the new
name began on 9 May 2025; the New York Stock Exchange ticker changed from
\texttt{GOLD} to \texttt{B}, while the Toronto Stock Exchange ticker remained
\texttt{ABX} \cite{barrick-name-2025}. The distinction matters whenever a
market-price series is stitched across the corporate action.

The corporate action does not interrupt the operating series. Team~4 extracted
its mine-level data from the quarterly reports published on the company's
investor website, under the Barrick Gold Corporation name until 2025 and under
the Barrick Mining Corporation name thereafter; both refer to the same issuer.
Team~4 adopts the process and cost definitions of the 2024 Annual Report and
checks its first forecast year against the 2025 report's guidance
\cite{team4,barrick-ar-2025}. The name and ticker change affects the market
series of Section~\ref{sec:lse-current-diagnostics} and the equity bridge,
but not the operating panel's continuity.

Barrick is a gold and copper exploration, development and mining company
\cite{barrick-about-2026}. Its original Team~4 operating layer is narrower:
eight gold mines, excluding copper and the remaining gold assets \cite{team4}.
Section~\ref{sec:copper-extension} adds Lumwana, Zald\'ivar and Jabal Sayid as
an explicit copper scenario, without rewriting the frozen gold experiment.
Neither perimeter represents every corporate asset or a reconciled claim on
the whole company. The case connects commodity prices, production and costs
to component margin and a conditional DCF, bringing the team's valuation
interests into its stochastic modelling framework.

\section{The company perimeter is a model input}
\label{sec:company-perimeter}

% Sources: Team 4 Sections 1.1.2--1.1.3, 2 and 3.1; Annual Report 2025,
% MD&A mine tables and Note 2; 2026Q1--Q2 mine statistics in Chapter 10.
A valuation must freeze, by asset and date, the ownership and reporting basis,
status and currency. These determine which volumes, costs, taxes, royalties,
cash flows, debt claims and shares enter its numerator and denominator.
Table~\ref{tab:company-operating-perimeter} records the gold perimeter.

\paragraph{Admitted assets.}
Bulyanhulu, Carlin, Cortez, Kibali, Loulo--Gounkoto, North Mara, Pueblo Viejo
and Turquoise Ridge form the base of the gold operating computations,
including Chapter~10's production-weighted cost master chain. The source
excluded copper and smaller gold operations \cite{team4}; this modelling
choice does not demonstrate that copper is economically negligible.
The eight mines produced 622 and 711 thousand attributable ounces in 2026Q1
and 2026Q2 against company gold totals 719 and 796 thousand ounces
(Chapter~10): 86.5\% and 89.3\% of attributable gold production. These are
gold-volume coverage ratios, not corporate-revenue or value coverage.

The frozen executable gold vector uses company totals of 719 and 796 koz,
and company-wide CoS of USD~1,922 and 1,993/oz, for its first two quarters.
From 2026Q3 it uses the eight-mine Team~4 forecast. Its first two actual
quarters therefore have a broader asset scope than its forecast perimeter.
This scope break remains explicit in the paired copper experiment; the
eight-mine coverage table does not reconcile it. A consistent mine-level
actual vector and matching aggregate costs are required before the baseline
can be read as a strict eight-mine valuation.

\begin{table}[H]
\centering
\caption{Team~4 gold perimeter: Barrick's interest and accounting treatment
at 31 December 2025 \cite{team4,barrick-ar-2025}.}
\label{tab:company-operating-perimeter}
\small
\renewcommand{\arraystretch}{1.12}
\begin{tabularx}{\textwidth}{@{}>{\raggedright\arraybackslash}p{2.15cm}%
>{\raggedright\arraybackslash}p{1.75cm}%
>{\raggedright\arraybackslash}p{3.85cm}%
>{\raggedright\arraybackslash}X@{}}
\toprule
\textbf{Mine} & \textbf{Location} & \textbf{Barrick interest and accounting}
& \textbf{Source-specific notes} \\
\midrule
Carlin & Nevada, USA & 61.5\% through Nevada Gold Mines; consolidated with
38.5\% non-controlling interest (Newmont). & Reference mine: 28 quarters,
2019Q1--2025Q4; simulated recovery in $[0.71,0.85]$. \\
Cortez & Nevada, USA & 61.5\% through Nevada Gold Mines; as Carlin. &
Winter seasonality in throughput. \\
Turquoise Ridge & Nevada, USA & 61.5\% through Nevada Gold Mines; as Carlin. &
Winter seasonality in throughput. \\
Pueblo Viejo & Dominican Republic & 60\%; consolidated with 40\%
non-controlling interest (Newmont). & Wet-season throughput seasonality. \\
Kibali & DR Congo & 45\%; equity-accounted joint venture (AngloGold Ashanti
45\%, DRC Government 10\%). & Outside consolidated revenue and CoS;
proportionate in attributable per-ounce metrics. \\
Loulo--Gounkoto & Mali & 80\%; consolidated with 20\% non-controlling
interest; deconsolidated 16 June, reconsolidated 16 December 2025. & Suspended
from 14 January; 29 thousand ounces in 2025; wet-season seasonality. \\
North Mara & Tanzania & 84\%; consolidated with 16\% non-controlling
interest (Government of Tanzania). & 25 quarters from separate reporting
in 2019Q4. \\
Bulyanhulu & Tanzania & 84\%; as North Mara. & -- \\
\bottomrule
\end{tabularx}
\end{table}

\paragraph{Reporting basis.}
The operating layer is attributable. Mine-level MD\&A tables report ore
processed, grade, recovery, ounces produced and sold and CoS on Barrick's
share: 61.5\% for the Nevada mines, 60\% Pueblo Viejo, 45\% Kibali and
84\% for the Tanzanian mines \cite{barrick-ar-2025}. Per-ounce CoS uses
attributable ounces sold. Ore and the ounces it generates share one basis;
Section~\ref{sec:company-operating-quantities}'s production identity is
dimensionally consistent. Grade and recovery are ratios. Ownership must not
be applied again to already attributable volumes.

Legal ownership does not always equal the economic share. Under the Tanzanian
agreement the government receives half of economic benefits through taxes,
royalties, clearing fees and cash distributions after capital recoupment;
earnings are recorded on Barrick's 84\% equity interest
\cite{barrick-ar-2025}. Volume reporting is suitable for production and cost;
the cash-flow split belongs to the tax and bridge layers.

\paragraph{Data window.}
The quarterly estimation window starts in 2019 and ends in 2025Q4 even where
longer histories exist \cite{team4}. Leading missing production-driver values
are removed mine by mine: Carlin contributes 28 quarters, North Mara 25 from
2019Q4. A common starting year therefore does not imply identical effective
samples for every driver. The 2026Q1--Q2 mine statistics replace the first
two forecast quarters as actuals; forecasts from 2026Q3 remain as originally
estimated (Chapter~10).

\section{Operating quantities supplied by the source programme}
\label{sec:company-operating-quantities}

% Sources: Team 4 Sections 1.1.3, 2.2--2.4, 3.1--3.2 and 4.1;
% Annual Report 2025, Production Costs and mine tables.
Production and CoS are forecast separately, mine by mine, on 20 quarters
from 2026Q1 to 2030Q4. Annual quantities sum four quarters. The price comes
from Team~8 in Chapter~9; Team~4's standalone demonstrator is excluded.

\paragraph{Production.}
Quarterly production of mine $i$ is built from three physical drivers,
\begin{equation}
\mathrm{Prod}_{i,t}=\frac{\kappa}{c}O_{i,t}G_{i,t}R_{i,t},
\qquad\kappa=0.91,\quad c=31.1035~\text{g/oz},
\label{eq:company-production-identity}
\end{equation}
where $O$ is ore processed in tonnes, $G$ processed grade in grams per tonne
and $R\in(0,1)$ recovery. Their product is recovered gold in grams; $c$
converts grams to troy ounces. The realisation factor $\kappa$ is the mean
mine ratio of realised to theoretical production, absorbing lags such as
stockpiling, heap leaching and autoclave residence \cite{team4}. The source
uses rounded $31.11$, lowering output approximately $0.02\%$; conversion
is from grams, not tonnes as stated in its text. Ore uses
SARIMA$(1,0,1)(1,0,0)_4$, grade AR(1) without seasonal terms, and recovery
uniform historical bounds widened by $\pm1\%$, truncated to $[0.50,1.00]$.
Production bands are 80\% Goodman product-variance intervals under
independence (Chapter~10).

\paragraph{Costs.}
CoS is USD per attributable ounce sold, excluding closure or care-and-
maintenance sites. The 2024 definition used by Team~4 includes site costs,
royalties, depreciation and community relations \cite{team4}; the 2025 report
also identifies mining and production taxes. Consolidated gold CoS in 2025
was USD~7,357 million: site costs 5,056; depreciation 1,588 (21.6\%);
royalties 540 (7.3\%); mining and production taxes 132; community relations
41 \cite{barrick-ar-2025}. The production-weighted log master chain is
ARIMA$(3,1,0)$ with drift, projected to each mine by
$\beta_i=0.67\,\sigma(r_{i,t})/\sigma(r_{\mathrm{global},t})+0.33$ and anchored
at 2025Q4. Cost bands are 95\%, versus production's 80\%.

\paragraph{Operating output.}
For quarterly gold price $P_t$ in USD/oz, the source passes
\begin{equation}
\mathrm{CM}_t=(P_t-\mathrm{CoS}_t)\mathrm{Prod}_t
\label{eq:company-component-margin}
\end{equation}
to valuation. Because CoS contains depreciation, this is after-depreciation
\emph{component margin}, not EBITDA despite the source label \cite{team4}.
Nevada Gold Mines' 2025 mine income USD~3,141 million plus depreciation 568
gives reported EBITDA 3,709 \cite{barrick-ar-2025}. The corporate gold cost
mix puts depreciation near one fifth of CoS. Nor is the margin reconciled
EBIT: G\&A, exploration, project expenditure and other items are absent.
CoS is per ounce \emph{sold}, while the identity multiplies ounces
\emph{produced}. Forecast royalties are fixed per ounce, although Barrick
attributes about USD~55/oz of the 2025 CoS increase to higher royalties from
higher realised gold prices \cite{barrick-ar-2025}. Chapter~11 must reconcile
these mismatches before the proxy becomes accounting cash flow.

\begin{table}[H]
\centering
\caption{Operating contract of the original Team~4 gold module.}
\label{tab:company-operating-contract}
\small
\renewcommand{\arraystretch}{1.12}
\begin{tabularx}{\textwidth}{@{}p{2.15cm}X X@{}}
\toprule
\textbf{Dimension} & \textbf{Team~4 specification} &
\textbf{Reconciliation still required} \\
\midrule
Perimeter & Eight gold mines, attributable; copper and other gold assets
absent from source module. & Separate copper extension; other assets,
per-share basis and North America IPO. \\
Frequency and horizon & Quarterly estimation 2019--2025; 20 quarters to
2030Q4; annual sums; Q1--Q2 2026 actuals. & No mine-life model supports a
terminal assumption beyond 2030. \\
Production & $O$ in t, $G$ in g/t, $R$ decimal, output oz; seasonal ore,
AR(1) grade, uniform recovery; 80\% Goodman bands. & Produced versus sold;
only point forecasts enter frozen valuation. \\
Costs & USD/oz CoS includes depreciation and royalties; log ARIMA$(3,1,0)$
with drift, $\beta$ scaling, Q4 2025 anchors; 95\% bands. & D\&A, royalties,
taxes, capex, G\&A, exploration and exceptional anchors; no AISC model. \\
Commodity price & Source only: GBM at USD~4,677/oz on 2 April 2026 using GLD
per-expiry volatility; excluded here. & Team~8 shapes at Q2 realised gold
USD~4,417/oz; no GLD-level conversion; realised basis, hedging and measures. \\
Operating output & After-depreciation component margin, called EBITDA in
source. & EBIT, taxes, working capital and capex bridge. \\
Dependence & Ore--grade $r\approx-0.45$, ore--recovery $r\approx-0.17$;
cost PC1 41.22\%, Bartlett $p=0.0186$; fixed operating vectors in valuation.
& Joint simulation, production response and royalty--price link. \\
\bottomrule
\end{tabularx}
\end{table}

\section{Financial and structural perimeter}
\label{sec:company-financial-perimeter}

The financial perimeter fixes statements, currency and scale for taxes,
reinvestment and capital costs, and debt, leases or equivalent claims,
non-controlling interests, non-operating assets and diluted shares. Model and
market claims must refer to the same security and date.

Barrick consolidates the Nevada, Pueblo Viejo, Loulo--Gounkoto and Tanzanian
operations and equity-accounts Kibali \cite{barrick-ar-2025}. Here flows
already represent Barrick's share. Deducting those same partners' interests
again would double count their exclusion. Consolidated cash and debt include
full subsidiary balances and exclude Kibali's; claims must be restated
consistently with attributable flows, including corporate-level claims,
not multiplied mechanically by one aggregate ownership fraction.

The 1,675.51 million diluted-share proxy remains unreconciled (Chapter~11).
Shares claim the whole company, whereas the gold baseline and gold-plus-
copper extension cover defined subsets. Per-share division allocates a
perimeter proxy; it does not yield full-company equity fair value. The copper
extension reports aggregate operating-value proxies without per-share targets.

Loulo--Gounkoto was suspended 14 January 2025, placed under provisional
administration and deconsolidated 16 June, then returned to Barrick control
16 December following the agreement announced 24 November. On 1 December
2025 the Board authorised preparation of an IPO of an entity holding its
Nevada Gold Mines and Pueblo Viejo interests plus Fourmile, expected by late
2026 with Barrick retaining control \cite{barrick-ar-2025}. External ownership
would change effective look-through interests in four modelled mines. The
December 2025 percentages cannot simply be carried beyond completion.
The 2025 Hemlo and Tongon sales lie outside the eight-mine vector.
The annual report alone does not fix IPO status on 2 September 2026.
The 28 April update and the 10 August Q2 release still describe a planned
minority-stake IPO expected by year end \cite{barrick-ipo-2026,barrick-q2-release-2026}.
The latter also announces Newmont's consent, an early contribution of excluded
properties to Nevada Gold Mines and a USD~1.95 billion top-up payment. Those
events require their own asset, timing and cash-bridge reconciliation; they
do not authorise a silent change to the frozen eight-mine forecast vector.
The 2026Q2 operating release predates the 2 September market snapshot, but its
reporting-period end, publication date and October retrieval remain distinct.

\gapbox{Freeze date, currency and scale; restate debt, cash, other claims and
non-operating assets consistently with attributable flows, without deducting
the same partners' interests twice; reconcile diluted shares; verify North
America IPO status at the valuation date and resulting look-through ownership;
document the security and corporate-action policy of the market comparison.}

\section{Risk boundary and permitted use}
\label{sec:company-risk-boundary}

The gold mines span Nevada, the Dominican Republic, Mali, Tanzania and the
DRC. Team~4's lag-four throughput term captures winter or wet-season
constraints. Structural events are not endogenous: the source treats the
Loulo--Gounkoto suspension as a mine-specific shock absorbed geographically
and over time, and simultaneous failures as negligible \cite{team4}.
This is an assumption, not an estimated geopolitical-risk bound.

Production diversification does not imply independent costs. Bartlett's test
rejects independence ($p=0.0186$); PC1 explains 41.22\% of quarterly cost
shocks, the first two components 62.79\% \cite{team4}. Shared cost shocks
motivate the master chain and remain systematic across the perimeter.

Cost paths start at 2025Q4. Loulo--Gounkoto's USD~4,151/oz CoS includes the
inventory fair-value increment on reconsolidation, versus cash costs
USD~1,448/oz. Its 2026 CoS guidance USD~2,860--3,140/oz still includes higher
purchase-price-allocation depreciation \cite{barrick-ar-2025}. Projecting
that anchor carries an exceptional effect forward; it needs a separate
normalisation check. This edit does not silently alter the frozen vector.

Reserve life is outside the model. Grade is stationary over five years under
feed-blending arguments; depletion, closure, high-grading, throttling and
care and maintenance are not endogenous \cite{team4}. Neither that model
nor flat extension of observed copper forwards supplies mine-life evidence
for perpetuity. The copper comparison uses a finite five-year primary
horizon and separate terminal sensitivities. Until accounting, structural
and measure bridges are admitted, the document defines a research contract
and its failure conditions, not a final target value.
